% This is samplepaper.tex, a sample chapter demonstrating the
% LLNCS macro package for Springer Computer Science proceedings;
% Version 2.20 of 2017/10/04
%
\documentclass[runningheads]{llncs}
%
\usepackage[numbers,sort&compress]{natbib}
\usepackage{hyperref}
\usepackage{amsmath,amsfonts}
\hypersetup{
    colorlinks=true,
    linkcolor=blue,
    filecolor=blue,  
    citecolor=blue,   
    urlcolor=blue,
}
\usepackage{graphicx}
\usepackage{subfigure}
\usepackage{xcolor}
\usepackage{multirow}

% Used for displaying a sample figure. If possible, figure files should
% be included in EPS format.
%
% If you use the hyperref package, please uncomment the following line
% to display URLs in blue roman font according to Springer's eBook style:
% \renewcommand\UrlFont{\color{blue}\rmfamily}

\newcommand{\xreal}{\boldsymbol{x}\textsubscript{real}}
\newcommand{\xfake}{\boldsymbol{x}\textsubscript{syn}}
\newcommand{\vx}{\boldsymbol{x}}
\newcommand{\vz}{\boldsymbol{z}}
\newcommand{\vf}{\boldsymbol{f}}
\newcommand{\vw}{\boldsymbol{w}}
\newcommand{\vI}{\boldsymbol{I}}
\newcommand{\vtheta}{\boldsymbol{\theta}}
\newcommand{\vbeta}{\boldsymbol{\beta}}
\newcommand{\TODO}[1]{\textcolor{red}{#1}}
\begin{document}
%
%\title{Survival Analysis of Cancer Patients from Portrait Photos Using StyleGAN Embedding}
\title{Deep Learning for Cancer Prognosis Prediction Using Portrait Photos by StyleGAN Embedding}
%
\titlerunning{Cancer Prognosis Prediction Using Portrait Photos by StyleGAN Embedding}
% If the paper title is too long for the running head, you can set
% an abbreviated paper title here
%
\author{Amr Hagag\inst{1,2}  \and Ahmed Gomaa\inst{1,3} \and Dominik Kornek\inst{1,3} \and
Andreas Maier\inst{2} \and\\ Rainer Fietkau\inst{1,3} \and Christoph Bert\inst{1,3} \and Florian Putz\inst{1,3} \and
Yixing Huang\inst{1,3}}
%
\authorrunning{A. Hagag et al.}
% First names are abbreviated in the running head.
% If there are more than two authors, 'et al.' is used.
%
\institute{Department of Radiation Oncology, University Hospital Erlangen, Friedrich-Alexander-Universit\"at Erlangen-N\"urnberg, Erlangen, Germany\\
\email{yixing.huang@uk-erlangen.de}\\
%\email{lncs@springer.com}\\
 \and
Pattern Recognition Lab, Friedrich-Alexander-Universit\"at Erlangen-N\"urnberg, Erlangen, Germany\\
\and
Comprehensive Cancer Center Erlangen-EMN (CCC ER-EMN), Erlangen, Germany
}
%
\maketitle              % typeset the header of the contribution
%
\begin{abstract}
Survival prediction for cancer patients is critical for optimal treatment selection and patient management. Current patient survival prediction methods typically extract survival information from patients' clinical record data or biological and imaging data. In practice, experienced clinicians can have a preliminary assessment of patients' health status based on patients' observable physical appearances, which are mainly facial features. However, such assessment is highly subjective. 
In this work, the efficacy of objectively capturing and using prognostic information contained in conventional portrait photographs using deep learning for survival predication purposes is investigated for the first time.
A pre-trained StyleGAN2 model is fine-tuned on a custom dataset of our cancer patients' photos to empower its generator with generative ability suitable for patients' photos. 
The StyleGAN2 is then used to embed the photographs to its highly expressive latent space.
Utilizing the state-of-the-art survival analysis models and based on StyleGAN's latent space photo embeddings, this approach achieved a C-index of $0.677$, which is notably higher than chance and evidencing the prognostic value embedded in simple 2D facial images. 
In addition, thanks to StyleGAN's interpretable latent space, our survival prediction model can be validated for relying on essential facial features, eliminating any biases from extraneous information like clothing or background. Moreover, a health attribute is obtained from regression coefficients, which has important potential value for patient care.
%which enables patients to edit their photos to appear healthier for better quality of life or provide a good prognosis of patients' facial appearance for better patient care.


\keywords{Survival prediction \and cancer \and StyleGAN \and deep learning \and latent space \and face prognosis \and explainable AI.}
\end{abstract}
%
%
%
\section{Introduction}
Survival analysis for cancer patients is critical for optimal treatment selection and patient management. Patients with very adverse prognosis may need aggressive management and acute intervention or may not benefit from tumor-directed treatment at all. Patients with good prognostic constellation similarly may benefit from specific treatment approaches like dose escalation and radical treatment of all sites in radiotherapy. So far, a vast class of survival prediction models have been investigated, ranging from conventional regression algorithms \cite{emura2012survival,kidd2018survival} to the latest radiomics \cite{van2017survival,chen2022mri} and deep learning methods \cite{wang2021cox,katzman2018deepsurv,kim2019deep,wankhede2022dynamic}.
With the successful applications in various fields including radiation oncology \cite{huang2022deep,weissmann2023deep,huang2023benchmarking,putz2023segment}, nowadays deep learning has become the main trend for patient survival prediction, because of its high representation power and ability to extract essential features from various types of data. Nevertheless, all the previous methods rely on data ranging from regular clinical records to commonly used medical imaging data (e.g., MRI and CT) as well as biological data (e.g., histologic and genetic information) \cite{vale2021long}. 
Within radiotherapy practice, the acquisition of patient photographs is a routine measure implemented for identity verification purposes.
Nevertheless, the potential utility in predicting patient survival and overall health status based on facial features derived from such widely available 2D portrait photos remains unexplored.

Facial features observable in 2D portrait photographs, like skin colorit and facial expression can reflect underlying pathophysiologic processes like muscle wasting, lowering of muscle tone or changes in microcirculation \cite{oken1982toxicity}. These features therefore hold considerable potential as indicators of a patient's underlying health status.
Many severe diseases will lead to evident facial changes, e.g., the most well-known diseases being genetic ones like the Down's syndrome \cite{gurovich2019identifying}. 
Assessing patients' overall health and prognosis based on the facial expression has a long-standing history and dates back at least to hippocrates \cite{withington2005hippocratic}.
Clinicians with extensive experience frequently utilize patients' physical appearance, as well as other related features, to estimate survival time and evaluate overall functional status. The Eastern Cooperative Oncology Group (ECOG) score is such a widely used metric for evaluating patient health status. This score is assigned by physicians and utilizes an ordinal 5-point scale to grade a patient's overall health \cite{oken1982toxicity}. 
Hui et al. \cite{hui2015diagnostic} were able to show that drooping of the nasolabial fold  was a highly predictive indicator of imminent death within a three-day timeframe.
Patient appearance has also been shown to be a factor in determining healthcare priority in emergency situations \cite{bagnis2020judging}. However, assessments made by human clinicians are inherently subjective and can therefore exhibit a high degree of intra- and inter-rater variance. In this work, we aim to achieve an objective assessment of cancer patients' prognosis based on facial features using deep learning.
 
Deep learning-based facial image analysis can reveal patient health status caused by various factors \cite{su2021deep}. It has been reported to be effective in detecting neurological disorders \cite{yolcu2017deep}, acromegaly \cite{kong2018automatic}, genetic disorders \cite{gurovich2019identifying}, and cancer \cite{liang2020identification}.
However, deep learning models can be biased \cite{wallis2022clever} and have poor generalizability. 
To ensure meaningful facial features are used, interpretability of the model is crucial. StyleGAN \cite{karras2020training,karras2020analyzing} is a widely used network for face image synthesis due to its visual fidelity, diversity, and high editability over its latent space. 
%Thus, the optimization of StyleGAN for generating latent vectors presents a promising avenue for survival analysis.
Essential prognostic features can potentially be embedded into the latent space as well and hence it presents a promising avenue for survival analysis. A notable advantage of leveraging the latent space for prognosis prediction is the ability to visually and qualitatively evaluate the associations between the features within the latent space and patient's prognosis.
 %such as genetic, neurological, psychiatric or medication-induced factors . It has been reported effective for monitoring neurological disorders \cite{yolcu2017deep}, detection of acromegaly \cite{kong2018automatic}, identifying genetic disorders \cite{gurovich2019identifying} and identifying the presence of cancer \cite{liang2020identification}. 
%However, deep learning models typically are black boxes. Many models with high accuracy in the training dataset actually have very poor generalizability as the models are actually using bias features \cite{wallis2022clever}. For survival prediction of cancer patients from portrait photos, many information in the photos have the risk to bias the model, for instance, the clothes and the objects in the background. Therefore, it is crucial to check whether the model is using meaningful facial features instead of such bias information, which requires a good interpretability of the deep learning model. The style-based generative adversarial network (StyleGAN) \cite{karras2020training,karras2020analyzing} is widely used for face image synthesis because of not only its prominent visual fidelity and diversity but also the high editability over its latent space \cite{abdal2019image2stylegan,abdal2020image2stylegan,roich2022pivotal,peterson2022deep}. Facial features can be well embedded into its $\mathcal{W}$ latent or extended $\mathcal{W}+$ latent space \cite{abdal2019image2stylegan} and hence photos can be easily edited to achieve desired styles such as face pose, orientation and expression. The effects of disentangled features in the latent space can be directly reflected in the altered synthetic photos, which actually provides good interpretability of each latent feature. In this work, we propose cancer patient survival prediction based on such StyleGAN latent features.

%This work contributes significantly to the field of cancer patient survival analysis using 2D facial photographs. %
Our contributions can be summarized as follows: a) To the best of our knowledge, this is the first attempt to use deep learning for cancer patient survival prediction based on 2D facial photographs;
b) We propose a model that utilizes the StyleGAN expressive latent space feature representations for survival prediction; 
c) To provide interpretability and validation to the model, a ``health" attribute is extracted to allow adjusting face photographs to either a healthier or more severe illness appearance. This manipulation has significant prognostic value for patient care and societal perception, underscoring its potential important clinical value.%importance for clinical applications.
% HYX: using a, b, c or 1), 2), 3) saves space and makes it straightforward to see how many points.

% The main contributions of this work lie in these aspects: a) To the best of our knowledge, our work is the first to apply deep learning for cancer patient survival prediction based on 2D facial photographs; b) A simple regression model is proposed for survival prediction utilizing the StyleGAN latent space features, where a 'health' attribute is extracted from the StyleGAN latent space for model interpretability; c) The health attribute allows photo editing into healthier or severer illness appearance, which has important potential value for patient care.


\begin{figure}[b]
\begin{minipage}[b]{0.16\linewidth}
\subfigure[$\xreal$]{
\includegraphics[width=\linewidth]{mask_original.png}
}
\end{minipage}
\begin{minipage}[b]{0.16\linewidth}
\subfigure[$\xfake$ \tiny{w/o FT}]{
\includegraphics[width=\linewidth]{mask_not-finetuned.png}
\label{subfig:noFT1}
}
\end{minipage}
\begin{minipage}[b]{0.16\linewidth}
\subfigure[$\xfake$ \tiny{with FT}]{
\includegraphics[width=\linewidth]{mask_fine-tuned.png}
\label{subfig:FT1}
}
\end{minipage}
\begin{minipage}[b]{0.16\linewidth}
\subfigure[$\xreal$ ]{
\includegraphics[width=\linewidth]{headcover_original.png}
}
\end{minipage}
\begin{minipage}[b]{0.16\linewidth}
\subfigure[$\xfake$ \tiny{w/o FT}]{
\includegraphics[width=\linewidth]{headcover_no_fine-tuned.png}
\label{subfig:noFT2}
}
\end{minipage}
\begin{minipage}[b]{0.16\linewidth}
\subfigure[$\xfake$ \tiny{with FT}]{
\includegraphics[width=\linewidth]{headcover_fine-tuned.png}
\label{subfig:FT2}
}
\end{minipage}
\caption{Exemplary face images and their corresponding synthetic images generated by StyleGAN without (w/o) and with fine-tuning (FT). The white rectangles are added for anonymization.}
\label{Fig:FineTuningEffect}
\end{figure}

\section{Material and Methods}
%\subsection{Data}
\textbf{2.1 Data}
\vspace{5pt}

The dataset utilized in this study comprises 16,657 2D portrait photographs of 13,523 individuals who have been diagnosed with cancer. Of this cohort, 52.5\% are male and 47.4\% are female. It is notable that 53.3\% of the dataset is right censored, indicating that the time of death for these individuals is either unknown or has not yet occurred. In addition, the survival times for the non-censored patients included in the dataset ranged from 2 to 4,923 days. To mitigate potential biases that may be present in the photographs due to factors such as clothing and background, an automated face recognition process utilizing dlib \cite{dlib09} is employed to align and crop all images to a standardized size of $512 \times 512$. The clinical information about the patients is retrieved, including international classification of diseases (ICD), gender, age, number of tumors, etc. Information acquired after the photo acquisitions is excluded.

%\subsection{Projecting Images to StyleGAN2 Latent Space}
\vspace{5pt}
\noindent\textbf{2.2 Embedding Images to StyleGAN2 Latent Space}
\vspace{5pt}

To accurately capture patients' facial features from a given photo, a generative model should be able to accurately reproduce these expressive features into a generated image.
A StyleGAN2 model pretrained on the FFHQ dataset \cite{karras2019style} is fine-tuned using adaptive discriminator augmentation (ADA) \cite{karras2020training} with our custom dataset of cancer patients' face images. This step is necessary to increase StyleGAN's generation ability for the special facial features in cancer patients and auxiliary medical equipment such as masks, breathing tubes and gauze. 
Fig.\,\ref{Fig:FineTuningEffect} shows the difference between synthesised images with and without fine-tuning. Without fine-tuning, the pretrained model fails to generate the mask and the bouffant cap in Fig.\,\ref{subfig:noFT1}  and Fig.\,\ref{subfig:noFT2}, respectively, while the fine-tuned model successfully re-synthesizes such structures.

% \subsection{Experimental Setup}
%\subsubsection{Dataset}

\noindent\textbf{Photo Embedding:} To project images from the image space to
the latent space, for each photo $\xreal$, an optimization process searches via gradient decent for a $512$-dimensional latent space vector $\vz$ if input to the StyleGAN2 generator, will produce a corresponding synthetic image $\xfake$ similar to the original image $\xreal$. This similarity is defined as Euclidean distance in the feature space of a pretrained VGG-16,
\begin{equation}
\vz = \arg \min_{\vz} \phi\left(\xreal, \xfake \right) = \arg \min_{\vz} \phi\left(\xreal, G(\vz) \right) ,
\end{equation}
where $G$ is the generator of the fine-tuned StyleGAN2 model, and $G(\vz)$ is the synthetic image $\xfake$, and $\phi$ is the loss function. 
% In such a way, a latent vector $\vz$ with a dimension of $N$ ($N=512$ as default) is learned for each image, representing the latent features of this image.
In such a way, each face image is associated with one optimized highly expressive $512$-dimensional latent vector.

%\subsection{Survival Analysis Models}
\vspace{5pt}
\noindent\textbf{2.3 Survival Analysis Models}
\vspace{5pt}

Survival analysis models are used to predict patients' survival by estimating the hazard function, which is the instantaneous rate at which the event of interest (e.g., death) occurs at a given time, and allows for the inclusion of multiple covariates to predict the risk of the event. 
%Given the known time-to-event of our patients' cohort, we propose to utilize StyleGAN photo embeddings $\vz$ to predict the continuous survival time.
The most widely used survial model in practice is the Cox proportional hazard (CoxPH) \cite{lin2002modeling} model based on conventional clinical features.
%\noindent\textbf{Cox proportional hazard (CoxPH) model.}
In the CoxPH model, the hazard function has the following form \cite{katzman2018deepsurv},
\begin{equation}
\lambda(t|\vf)=\lambda_0(t)\cdot e^{h(\vf)}=\lambda_0(t)\cdot e^{\vw^\top\vf},
\end{equation}
where $\lambda_0(t)$ is a baseline hazard function, $\vf$ is the input feature vector, $h(\vf)$ is the log-risk function, and $\vw$ is the coefficient vector to optimize. 

%\noindent\textbf{DeepSurv.}
To learn a more complex, non-linear relationship between the features for the hazard estimation, deep learning networks have been proposed, with DeepSurv \cite{zhu2016deep,katzman2018deepsurv} being the state-of-the-art. The original DeepSurv method applies a multi-layer perceptron (MLP) to estimate the log-risk function $h(\vf)$ of the CoxPH model. To extract features from medical images, the MLP network can be replaced by a convolutional neural network (CNN) \cite{zhu2016deep,gurovich2019identifying}. In this work, a ResNet-18 and a VGG-16 network are investigated as examples of the CNN with image data, which are trained in an end-to-end approach. As CNN-extracted features are highly abstract, in this work, we propose to integrate the StyleGAN network to the DeepSurv framwork for the prediction of continuous survival time, i.e., utilizing the expressive StyleGAN photo embedding $\vz$ as the input features of the CoxPH and DeepSurv models. Our method uses a two-step training with the benefit of good explainability. The first step provides a comprehensive set of features, thanks to the well-established StyleGAN latent space. The second step selects the features from this set, which are very important for survival prognosis. Such features can be later visualized directly by the existing StyleGAN generator.

%DeepSurv \cite{katzman2018deepsurv} combines the CoxPH model with a neural network to learn a complex, non-linear relationship between the covariates. This allows the model in our case to capture more complex relationships and interactions between the dimensions of the latent vector, which may be missed by conventional CoxPH model. Based on the input data, DeepSurv can utilize a multi-layer perceptron (MLP) or a convolutional neural network (CNN) to estimate the log-risk function $h(\vf)$ of the CoxPH model. In this work, a ResNet-18 and a VGG-16 network are investigated as examples of the CNN with image data.
%The objective function is \cite{katzman2018deepsurv}, \TODO{(equation can be removed to save space)}
% \begin{equation}
% l(\theta) = -\frac{1}{N_{E=1}}\sum_{i:E_i=1}\left(\hat{h}_{\vtheta}(\vf_i)-\log \sum_{j\in \mathcal{R}(T_i)}e^{\hat{h}_{\vtheta}(\vf_j)} \right) + \lambda||\vtheta||_2^2,
% \end{equation}
%where $\vtheta$ is the network parameter set and $\hat{h}_{\vtheta}$ is the network output, $\lambda$ is a regularization parameter.

% \textbf{DeepConvSurv.}
% The MLP network in the DeepSurv is replaced by a convolutional neural network (CNN) to extract features from photo images to estimate the log-risk function $h(\vf)$. The function of the CNN is similar to that in \cite{gurovich2019identifying} but the network is trained to fit $h(\vf)$ for the CoxPH model. In this work, a ResNet and a VGG-16 network are investigated as examples of the CNN. 

% \subsection{Survival prediction using StyleGAN2 latent features}
% A StyleGAN2 model pretrained on the FFHQ dataset \cite{karras2019style} is fine-tuned using adaptive discriminator augmentation (ADA) \cite{karras2020training} with our custom dataset of cancer patients' face images. This step is necessary to increase the StyleGAN's generation ability for the special facial features in cancer patients and auxiliary medical equipment such as masks, breathing tubes and gauze, both of which typically do not occur frequently in the FFHQ dataset (healthy population). 

% For each cancer patient photo $\xreal$, a latent space code $\vz$ is trained to minimize the loss between the original image $\xreal$ and its synthetic image $\xfake$,
% \begin{equation}
% \vz = \arg \min_{\vz} \phi\left(\xreal, \xfake \right) = \arg \min_{\vz} \phi\left(\xreal, G(\vz) \right) ,
% \end{equation}
% where $G$ is the generator of the fine-tuned StyleGAN model, and $G(\vz)$ is the synthetic image $\xfake$, and $\phi$ is  the VGG-16 perceptual loss. In such a way, a latent vector $\vz$ with a dimension of $N$ ($N=512$ as default) is learned for each image, representing the latent features of this image. 
% We propose to construct survival analysis models to predict the continuous survival time utilizing the photo embeddings $\vz$. 

\begin{figure}[t]
\includegraphics[width=1\linewidth]{MICCAI_Diagram.pdf}
\caption{The workflow of our proposed facial prognosis method: (a) Optimization of $\vz$ using StyleGAN generator fine-tuned on custom cancer patient data; (b) CoxPH regression to get $\vw$ for health attribute construction and latent space manipulation; (c) DeepSurv for survival prediction; (d) Generation of new appearance.}
\end{figure}

%\subsection{Latent Space Manipulation}
\vspace{5pt}
\noindent\textbf{2.4 Latent Space Manipulation}
\vspace{5pt}

In this work we propose two tests to visually and qualitatively evaluate the semantic meaning of our photo embedding as well as the prognostic value it holds for the task of survival prediction. Each of the tests can be conducted by simple manipulations of the latent vector $\vz$. 

\noindent\textbf{Synergetic Effect of Joint Features:} Leveraging the linearity of the CoxPH regression model, we fitted the $512$-dimensional feature vectors $\vz$ to the known survival data,
% Different regression models can be applied for survival prediction. The simplest regression model is the linear regression model,
\begin{equation}
h(\vz) =\sum_{i=0}^{N-1}\vw_i \vz_i,
\end{equation}
where $\vw$ is the vector of linear coefficients, and $h(\vz)$ is the log-risk function which is fed to the CoxPH model. By manipulating $\vz$ in all the dimensions jointly with a multiplier of $\vw$,
\begin{equation}
\hat{\vz} = \vz + \beta \vw,
\label{eqn:jointFeature}
\end{equation}
the dimensions in $\vz$ with high importance are modified to a large degree, while those with low importance stay almost unchanged. By changing $\beta$, the StyleGAN generator is able to generate new appearances of this patient from $\hat{\vz}$, with a healthier or more severe illness appearance. This vector $\vw$ serves as a health attribute.
% \noindent\textbf{Effect of Each Individual Feature:} Each dimension of the latent code $\vz$ can be manipulated one by one. The $i$-th dimension can be changed as follows,
% \begin{equation}
% \hat{\vz} = \vz + \beta \vI_i,
% \end{equation}
% where $\vI_i$ is a zero vector except that its $i$-th element equals 1. The magnitude of each weight $\vw_i$ straightforwardly reflects the importance of its corresponding latent feature for survival prediction. If important attributes manipulate unrelated appearance, e.g. clothes and background, instead of the patient facial appearance, the survival model is very likely biased.

\noindent\textbf{Age Manipulation:} To further validate the quality of our facial features vectors and ensure that our predictive model is focusing on complex survival related aspects and not only a single aspect like the age, we use a simple linear regression to map the feature vectors $\vz$ to the known age of the patients. Similar to (Eqn.\,(\ref{eqn:jointFeature})), $\vw$ here would serve as a simple age attribute.


%\subsection{Experimental parameters}
\vspace{5pt}
\noindent\textbf{2.5 Experimental Parameters}
\vspace{5pt}

\noindent\textbf{StyleGAN Fine-Tuning:} The StyleGAN version $2$ was utilized and fine-tuned using transfer learning from the FFHQ dataset of human faces, which had a resolution of $512 \times 512$. The fine-tuning process took $88$ hours and resulted in a final FID score of $5.1$ after $3200$k images. The fine-tuning was performed on a single NVIDIA Quadro RTX $8000$ $48$GB GPU.

\noindent\textbf{Iterative Embedding Optimization:}\ The optimization process was initialized with the latent vector $\overline{\vz}$ of the average face and performed $800$ gradient descent steps for the optimization using Adam optimizer with an initial learning rate of $0.01$, $\beta_1 = 0.9, \beta_2 = 0.999, \text{ and } \epsilon = 1e^{-8}$

\noindent\textbf{DeepSurv:} In this work, the DeepSurv model was trained with a train-test split of 80\% for training and validation and 20\% for testing using a batch size of $128$ and with the Adam optimizer and a learning rate of $0.001$. The network architecture consisted of a MLP with $6$ layers with batch normalization and dropout of $0.4$ to enhance the model's generalization and prevent overfitting. Furthermore, the Kaplan-Meier method was employed to estimate the censoring distribution, which is crucial for survival analysis tasks.

\noindent\textbf{Data Anonymization Disclaimer:} The data presented in this work has been anonymized to protect the privacy of individuals. Any personal information has been removed or de-identified, and the data cannot be used to identify any specific individuals. Facial images are either synthetic, anonymized, or projections that have been reviewed by clinical experts to ensure that no identification of individuals is possible.

% In this work, the following two ways are used to check whether bias information is used for survival prediction.


\begin{table}[b]
\caption{Survival prediction results of different methods using DeepSurv (DS) and CoxPH models. C-index and Brier score (B-Score) is caculated for all models.}
\label{Tab:accuracy}
\centering
\begin{tabular}{l|l r r||l| l r r}
\hline
\textbf{Input data}& \textbf{Method} & \textbf{C-index} & \textbf{B-Score} & \textbf{Input data} & \textbf{Method} & \textbf{C-index} & \textbf{B-Score} \\
\hline
Embeddings & DS+MLP &0.677 & 0.130 & Clinic. Data & CoxPH & 0.690 & 0.151\\ 
 
% \cline{2-4}\cline{6-8}
 Clinic. Data&DS+MLP  & 0.729 & 0.168 &Photos& DS+ResNet& 0.510 & 0.471\\
% \cline{1-8}
Combined& DS+MLP& 0.787 & 0.137 &Photos &DS+VGG & 0.470 & 0.471 \\
\hline
\end{tabular}
\end{table}


% \begin{table}[b]
% \caption{The survival prediction results of different methods. \TODO{TODO} }
% \label{Tab:accuracy}
% \centering
% \begin{tabular}{l|l|c|c||l|l|c|c}
% \hline
% \textbf{Input data}& Method & C-index & Brier-Score & Input data & Method & C-index & Brier-Score \\
% \hline
% \multirow{2}{*}{Clinical data} &CoxPH & 0.690 & 0.119 & \multirow{2}{*}{Photos} & DS+ResNet& 0.510 & 0.0\\
% % \cline{2-4}\cline{6-8}
%  &DS+MLP  & 0.729 & 0.168 & &DS+VGG & 0.470 & 0.0 \\
% % \cline{1-8}
% Combined& DS+MLP& 0.759 & 0.159 & SG Embedding&DS+MLP &0.677 & 0.130 \\
% % \hline
% \end{tabular}
% \end{table}

% \textbf{Experimental parameters}
% \TODO{Please report the training parameters here. For example, the number of epochs for fine-tuning the StyleGAN model. What is the $\lambda$ value for DeepSurf. What are the optimizers and epochs for the CNN models. And so on.} The StyleGAN version 2 is used in this work.
% \begin{figure}[t]
% \centering
% \includegraphics[width=1 \linewidth]{coefs_plot.png}
% \caption{The coefficients vector of the 512 latent features.}
% \label{Fig:WCoefficients}
% \end{figure}

\section{Results and Discussion}
%\subsection{Survival prediction results}
\textbf{3.1 Survival Prediction Results}
% \vspace{5pt}

The survival prediction results are displayed in Table\,\ref{Tab:accuracy}. 
Our approach for continuous survival time prediction using DeepSurv based on StyleGAN photo embeddings achieved a mean concordance index (C-index) of $0.677$ and a mean Brier Score of $0.130$. To put this into perspective, the model's survival predictions are evidently higher than chance. 
This reflects that the model uses key survival features from relevant prognostic facial attributes embedded in the latent vectors. As also shown in Table\,\ref{Tab:accuracy}, the DeepSurv model using patient's portraits fed directly through a CNN achieved a C-index of $0.47$ and $0.51$ using the ResNet-18 and the VGG-16 networks respectively.
In this case, the CNNs failed to extract the needed prognostic information directly from the portrait photos. 

As also shown in Table\,\ref{Tab:accuracy}, survival prediction models based on conventional clinical features achieved a C-index of 0.691 and 0.729 using CoxPH and DeepSurv models respectively, which is comparable to what was achieved by the embeddings of simple 2D portraits. 
Combining clinical features with photo embeddings boosted the C-index to 0.787 highlighting again the complementary prognostic value of StyleGAN's embeddings. It is worth mentioning that the combination was done by early fusion, in which the inputs were simply concatenated before feeding them into the network. A more advanced fusion method could potentially produce better results, which is our future work.

\begin{figure}[t]
\centering
\begin{minipage}[b]{0.16\linewidth}
\centering
$\beta = -10$
\end{minipage}
\begin{minipage}[b]{0.16\linewidth}
\centering
$\beta = -5$
\end{minipage}
\begin{minipage}[b]{0.16\linewidth}
\centering
$\xfake, \beta = 0$
\end{minipage}
\begin{minipage}[b]{0.16\linewidth}
\centering
$\beta = 5$
\end{minipage}
\begin{minipage}[b]{0.16\linewidth}
\centering
$\beta = 10$ 
\end{minipage}
\begin{minipage}[b]{0.16\linewidth}
\centering
$\beta = 20$ 
\end{minipage}

\begin{minipage}[b]{0.16\linewidth}
\includegraphics[width=\linewidth]{health_Beta_-10.0.png}
\end{minipage}
\begin{minipage}[b]{0.16\linewidth}
\includegraphics[width=\linewidth]{health_Beta_-5.0.png}
\end{minipage}
\begin{minipage}[b]{0.16\linewidth}
\includegraphics[width=\linewidth]{health_Beta_0.0.png}
\end{minipage}
\begin{minipage}[b]{0.16\linewidth}
\includegraphics[width=\linewidth]{health_Beta_5.0.png}
\end{minipage}
\begin{minipage}[b]{0.16\linewidth}
\includegraphics[width=\linewidth]{health_Beta_10.0.png}
\end{minipage}
\begin{minipage}[b]{0.16\linewidth}
\includegraphics[width=\linewidth]{health_Beta_20.0.png}
\end{minipage}

\begin{minipage}[b]{0.16\linewidth}
\includegraphics[width=\linewidth]{5_health_beta_-10.png}
\end{minipage}
\begin{minipage}[b]{0.16\linewidth}
\includegraphics[width=\linewidth]{5_health_beta_-5.png}
\end{minipage}
\begin{minipage}[b]{0.16\linewidth}
\includegraphics[width=\linewidth]{5_health_beta_0.png}
\end{minipage}
\begin{minipage}[b]{0.16\linewidth}
\includegraphics[width=\linewidth]{5_health_beta_5.png}
\end{minipage}
\begin{minipage}[b]{0.16\linewidth}
\includegraphics[width=\linewidth]{5_health_beta_10.png}
\end{minipage}
\begin{minipage}[b]{0.16\linewidth}
\includegraphics[width=\linewidth]{5_health_beta_-20.png}
\end{minipage}

\begin{minipage}[b]{0.16\linewidth}
\includegraphics[width=\linewidth]{2_health_Beta_-10.0.png}
\end{minipage}
\begin{minipage}[b]{0.16\linewidth}
\includegraphics[width=\linewidth]{2_health_Beta_-5.0.png}
\end{minipage}
\begin{minipage}[b]{0.16\linewidth}
\includegraphics[width=\linewidth]{2_health_Beta_0.0.png}
\end{minipage}
\begin{minipage}[b]{0.16\linewidth}
\includegraphics[width=\linewidth]{2_health_Beta_5.0.png}
\end{minipage}
\begin{minipage}[b]{0.16\linewidth}
\includegraphics[width=\linewidth]{2_health_Beta_10.0.png}
\end{minipage}
\begin{minipage}[b]{0.16\linewidth}
\includegraphics[width=\linewidth]{2_health_Beta_20.0.png}
\end{minipage}
\caption{Manipulating the health attribute with different $\beta$ values for regenerated patients' images. The synthesized image without manipulation is at $\beta=0$ (third column). These are projections that do not allow identification of real persons}
\label{Fig:JointFeatures}
\end{figure}
%\subsection{Model interpretability with ``health" attribute}
\vspace{5pt}
\noindent\textbf{3.2 Model Interpretability with ``Health" Attribute}
\vspace{5pt}

% The coefficient vector $\vw$ of the CoxPH linear regression model using StyleGAN latent features is displayed in Fig.\,\ref{Fig:WCoefficients}, where certain dimensions have obviously more importance than the others for survival prediction task. 
Fig.\,\ref{Fig:JointFeatures} demonstrates the synergetic effect of manipulating patient's images in the latent space using the health attribute
 (Eqn.\,(\ref{eqn:jointFeature})). In Fig.\,\ref{Fig:JointFeatures} the appearances of the patients are controlled by $\beta$. With a negative value for $\beta$, the patients appear to be healthier with a reddish skin colorit, full cheeks and a smiling expression with a strong tone of the facial muscles.
 In contrast, with a positive value for $\beta$, there is loss of masticatory muscles and subcutaneous fat with a pale skin colorit. Moreover, the tone of the facial muscles is reduced with drooping of the nasolabial fold, an overall more relaxed appearance of facial features and an increased ``tippiness" of the nose. These features are consistent with long-standing clinical descriptions of prognostically relevant changes in facial appearance \cite{hui2015diagnostic, withington2005hippocratic}. Most importantly, the change is mainly located in the head region instead of their clothes or the background. The model is also able to reflect patients' mental status to some degree by extracting their facial expressions.
 
 This editability of patients' photos to have either a healthier appearance or prognosis of a more severe illness has important value for patient care. For cancer patients, the quality of life is as important as the prevention, treatment and management of their illness. The quality of life relies on not only physical but also mental well-being \cite{rankin2003perceived,de2017quality}. Editing their photos to have a healthier appearance and sharing such optimized photos to their social networks can convey a more positive and optimistic message. For healthcare professionals, this could enable illustrating a patients' appearance in a healthy state or a state of illness progression, which could improve empathy and caretaking.
 
 An observation that could not be neglected however is the similarity between progressing in age and worse prognostic features, as age is a common risk indicator. While one can assume an underlying correlation between age and general health related facial features, we addressed a potential bias of the model to only rely on age features for the prognosis prediction task. 
 At the top row of Fig.\,\ref{Fig:JointFeatures}, it is clearly observable that when editing the patient's image based on the health attribute, distinct facial features are affected compared to when editing based on the age attribute, as shown at the top row of Fig.\,\ref{Fig:IndividualImportanceFeature}.
 
Based on each feature's importance for the task of survival prediction, a evident change is observed in the manipulated image. The images in the second row of Fig.\,\ref{Fig:IndividualImportanceFeature} that changing random features in the latent space would barely change features with any prognostic value in a meaningful way. This demonstrates that, as in Fig.\,\ref{Fig:JointFeatures}, the StyleGAN regression model is indeed using meaningful facial features instead of bias or random information for survival prediction.
 
 \begin{figure}[t]
\centering
\begin{minipage}[b]{0.16\linewidth}
\centering
$\beta = -10$
\end{minipage}
\begin{minipage}[b]{0.16\linewidth}
\centering
$\beta = -5$
\end{minipage}
\begin{minipage}[b]{0.16\linewidth}
\centering
$\xfake, \beta = 0$ %$\xreal$
\end{minipage}
\begin{minipage}[b]{0.16\linewidth}
\centering
$\beta = 5$
\end{minipage}
\begin{minipage}[b]{0.16\linewidth}
\centering
$\beta = 10$ 
\end{minipage}
\begin{minipage}[b]{0.16\linewidth}
\centering
$\beta = 20$ 
\end{minipage}

\begin{minipage}[b]{0.16\linewidth}
\includegraphics[width=\linewidth]{age_Beta_-3.0.png}
\end{minipage}
\begin{minipage}[b]{0.16\linewidth}
\includegraphics[width=\linewidth]{age_Beta_-1.5.png}
\end{minipage}
\begin{minipage}[b]{0.16\linewidth}
\includegraphics[width=\linewidth]{health_Beta_0.0.png}
\end{minipage}
\begin{minipage}[b]{0.16\linewidth}
\includegraphics[width=\linewidth]{age_Beta_1.0.png}
\end{minipage}
\begin{minipage}[b]{0.16\linewidth}
\includegraphics[width=\linewidth]{age_Beta_3.0.png}
\end{minipage}
\begin{minipage}[b]{0.16\linewidth}
\includegraphics[width=\linewidth]{age_Beta_6.0.png}
\end{minipage}

\begin{minipage}[b]{0.16\linewidth}
\includegraphics[width=\linewidth]{face_D132_-5.png}
\end{minipage}
\begin{minipage}[b]{0.16\linewidth}
\includegraphics[width=\linewidth]{face_D132_-2.6.png}
\end{minipage}
\begin{minipage}[b]{0.16\linewidth}
\includegraphics[width=\linewidth]{face_D132_0.png}
\end{minipage}
\begin{minipage}[b]{0.16\linewidth}
\includegraphics[width=\linewidth]{face_D132_1.7.png}
\end{minipage}
\begin{minipage}[b]{0.16\linewidth}
\includegraphics[width=\linewidth]{face_D132_2.3.png}
\end{minipage}
\begin{minipage}[b]{0.16\linewidth}
\includegraphics[width=\linewidth]{face_D132_3.png}
\end{minipage}

% \begin{minipage}[b]{0.16\linewidth}
% \includegraphics[width=\linewidth]{beard_432_Beta_-6.0.png}
% \end{minipage}
% \begin{minipage}[b]{0.16\linewidth}
% \includegraphics[width=\linewidth]{beard_432_Beta_-2.0.png}
% \end{minipage}
% \begin{minipage}[b]{0.16\linewidth}
% \includegraphics[width=\linewidth]{beard_432_Beta_0.0.png}
% \end{minipage}
% \begin{minipage}[b]{0.16\linewidth}
% \includegraphics[width=\linewidth]{beard_432_Beta_2.0.png}
% \end{minipage}
% \begin{minipage}[b]{0.16\linewidth}
% \includegraphics[width=\linewidth]{beard_432_Beta_5.0.png}
% \end{minipage}
% \begin{minipage}[b]{0.16\linewidth}
% \includegraphics[width=\linewidth]{beard_432_Beta_8.0.png}
% \end{minipage}
\caption{Exemplary patients' images generated by manipulating the age attribute (top row), and a single feature at random dimension $132$ (bottom row). These are projections that do not allow identification of real persons.}
\label{Fig:IndividualImportanceFeature}
\end{figure}

% This work demonstrates that our StyleGAN-based survival prediction model indeed extracts key survival features from relevant prognostic facial attributes such as the skin pigmentation and muscle tone instead of bias information in the form of an embedding vector.
% Furthermore, the health attribute empowers the StyleGAN with a good edibility of patients' photos to have either healthier or a more severe illness appearance. This has important value for both patients and clinicians. For cancer patients, the quality of life is as important as the prevention, treatment and management of their illness. The quality of life relies on not only physical but also mental well-being (e.g. social networking) \cite{rankin2003perceived,de2017quality}. Editing their photos to have a healthier appearance will convey a more positive message and make it easier for the patients to share their recent statuses to their friends. For clinicians, a good prognosis of patients' facial appearance is also very important for patient care. For example, the foreseeable muscle loss of patients' faces is beneficial for treatment decisions of head and neck cancer and the obvious deviation of patients' facial appearance to clinicians' expectation will help clinicians notice the side effects of the treatment earlier (e.g., swelling after radiation therapy).
We acknowledge a potential bias w.r.t ethnicity. However, unlike MRI and X-ray data, anonymization of facial photos is not possible, which restricts multicenter collaboration and leads to a lack of publicly available datasets. Additionally, our test split includes $>$3k samples, which is still larger than the training splits found in most datasets in the literature.

Using a dataset with only cancer patients was intentional, as the model is specifically designed for survival prediction of cancer patients, not for the general population. Nevertheless, the survival time in our dataset ranges from 3 days to 13 years, covering a diverse range of health statuses. Compared with the short-survival patients, long-survival patients are relatively healthier. Many of them have a long survival time just close to people without cancer. Therefore, health features can still be extracted.

In this preliminary work, we intentionally decided not to split the dataset into multiple subgroups, as the clinical literature indicated that facial prognostic features are universal and transcend cancer types. Moreover, splitting solely by tumor type also would not have addressed additional clinical dimensions like disease stage, specific molecular subtype etc. and was expected to greatly reduce the generalizability of the model. Therefore, we included images that were captured when the patient was admitted for radiotherapy at various disease stages. For future work, the performance of our method on survival analysis of different cancer groups will be investigated.


Currently, the well-established cancer prognosis standard is the AJCC TNM staging \cite{amin2017eighth}, which incorporates tumor, lymph nodes, and metastasis as the primary cues. Our model provides an additional tool for prognosis estimation, which is less influenced by human’s wishful thinking, complementing traditional clinical variables. Consequently, we showed a significant improvement in prognostic accuracy over relying solely on clinical variables, thereby supporting more effective oncologic decision-making.

\section{Conclusion}
This work investigates on prognosis prediction of cancer patients using portrait photos. Our work demonstrates that the latent features extracted by StyleGAN contain highly relevant prognostic information, and an effective regression model can be optimized for survival prediction. Such a model using facial features have good interpretability and can achieve comparable survival prediction accuracy to a state-of-art model using clinical record data. The constructed health attribute allows photo editing into healthier appearance or prognosis of a more severe illness, which has important potential value for patient care. %Among other applications, this could enable illustrating for health care professionals a patients' real appearance in a healthy state without signs of illness, which could improve empathy and caretaking.

%\appendix
%
%
%\section{StyleGAN Fine-Tuning}
%
%We build upon the official PyTorch implementation of StyleGAN2-ADA by Karras et al. \cite{Karras2020ada}, from which we inherit most of the training details.
%A pre-trained StyleGAN2 model was utilized and fine-tuned using transfer learning from the FFHQ dataset of human faces \cite{karras2019style}, which had a resolution of $512 \times 512$. The images in the custom dataset of cancer patients, on which the model was fine-tuned, were preprocessed and cropped to match the preprocessing of the face images in the FFHQ dataset. Fig.\,\ref{fig:fid} shows the progress of the training with relation to time and number of kimages.
%Fig.\,\ref{fig:fakes} demonstrates exemplary fake images generated by the fine-tuned generator.
%
%% \begin{figure}[h]
%% \includegraphics[width=0.7\linewidth]{FID_kimg.png}
%% \caption{Progress of Fréchet inception distance as a function of kimage during the model fine tuning}
%% \end{figure}
%
%% \begin{figure}[h]
%% \includegraphics[width=0.7\linewidth]{FID_hours.png}
%% \caption{Progress of Fréchet inception distance as a function of hours during the model fine tuning}
%% \end{figure}
%
%
%\begin{figure}[h]
%\begin{minipage}[b]{0.50\linewidth}
%\includegraphics[width=\linewidth]{FID_kimg.png}
%\end{minipage}
%\begin{minipage}[b]{0.5\linewidth}
%\includegraphics[width=\linewidth]{FID_hours.png}
%\end{minipage}
%\caption{Progress of Fréchet inception distance (FID) during the model fine tuning as a function of hours (on the right) and as a function of kimage (on the left)}
%\label{fig:fid}
%\end{figure}
%
%\begin{figure}[h]
%\includegraphics[width=1\linewidth]{fakes003200_cropped_resized_1.png}
%\caption{Exemplary fakes generated at resolution of $512 \times 512$ by the fine-tuned StyleGAN model after $3200$k images}
%\label{fig:fakes}
%\end{figure}
%
%
%\section{``Health" Attribute Manipulation}
%%Fig.\,\ref{fig:manipulations_1} and 
%Fig.\,\ref{fig:manipulations_2} show exemplary images for manipulating the health attribute with different $\beta$ values for regenerated patients' images.
%
%\begin{figure}[h]
%\centering
%\begin{minipage}[b]{0.134\linewidth}
%\centering
%$\beta = -20$
%\end{minipage}
%\begin{minipage}[b]{0.134\linewidth}
%\centering
%$\beta = -10$
%\end{minipage}
%\begin{minipage}[b]{0.134\linewidth}
%\centering
%$\beta = -5$
%\end{minipage}
%\begin{minipage}[b]{0.134\linewidth}
%\centering
%$\xfake, \beta = 0$
%\end{minipage}
%\begin{minipage}[b]{0.134\linewidth}
%\centering
%$\beta = 5$
%\end{minipage}
%\begin{minipage}[b]{0.134\linewidth}
%\centering
%$\beta = 10$ 
%\end{minipage}
%\begin{minipage}[b]{0.134\linewidth}
%\centering
%$\beta = 20$ 
%\end{minipage}
%
%\begin{minipage}[b]{0.134\linewidth}
%\includegraphics[width=\linewidth]{1001632229_beta_-20.0.png}
%\end{minipage}
%\begin{minipage}[b]{0.134\linewidth}
%\includegraphics[width=\linewidth]{1001632229_beta_-10.0.png}
%\end{minipage}
%\begin{minipage}[b]{0.134\linewidth}
%\includegraphics[width=\linewidth]{1001632229_beta_-5.0.png}
%\end{minipage}
%\begin{minipage}[b]{0.134\linewidth}
%\includegraphics[width=\linewidth]{1001632229_beta_0.0.png}
%\end{minipage}
%\begin{minipage}[b]{0.134\linewidth}
%\includegraphics[width=\linewidth]{1001632229_beta_5.0.png}
%\end{minipage}
%\begin{minipage}[b]{0.134\linewidth}
%\includegraphics[width=\linewidth]{1001632229_beta_10.0.png}
%\end{minipage}
%\begin{minipage}[b]{0.134\linewidth}
%\includegraphics[width=\linewidth]{1001632229_beta_20.0.png}
%\end{minipage}
%
%\begin{minipage}[b]{0.134\linewidth}
%\includegraphics[width=\linewidth]{1002102777_beta_-20.0.png}
%\end{minipage}
%\begin{minipage}[b]{0.134\linewidth}
%\includegraphics[width=\linewidth]{1002102777_beta_-10.0.png}
%\end{minipage}
%\begin{minipage}[b]{0.134\linewidth}
%\includegraphics[width=\linewidth]{1002102777_beta_-5.0.png}
%\end{minipage}
%\begin{minipage}[b]{0.134\linewidth}
%\includegraphics[width=\linewidth]{1002102777_beta_0.0.png}
%\end{minipage}
%\begin{minipage}[b]{0.134\linewidth}
%\includegraphics[width=\linewidth]{1002102777_beta_5.0.png}
%\end{minipage}
%\begin{minipage}[b]{0.134\linewidth}
%\includegraphics[width=\linewidth]{1002102777_beta_10.0.png}
%\end{minipage}
%\begin{minipage}[b]{0.134\linewidth}
%\includegraphics[width=\linewidth]{1002102777_beta_20.0.png}
%\end{minipage}
%
%
%\begin{minipage}[b]{0.134\linewidth}
%\includegraphics[width=\linewidth]{1002212612_beta_-20.0.png}
%\end{minipage}
%\begin{minipage}[b]{0.134\linewidth}
%\includegraphics[width=\linewidth]{1002212612_beta_-10.0.png}
%\end{minipage}
%\begin{minipage}[b]{0.134\linewidth}
%\includegraphics[width=\linewidth]{1002212612_beta_-5.0.png}
%\end{minipage}
%\begin{minipage}[b]{0.134\linewidth}
%\includegraphics[width=\linewidth]{1002212612_beta_0.0.png}
%\end{minipage}
%\begin{minipage}[b]{0.134\linewidth}
%\includegraphics[width=\linewidth]{1002212612_beta_5.0.png}
%\end{minipage}
%\begin{minipage}[b]{0.134\linewidth}
%\includegraphics[width=\linewidth]{1002212612_beta_10.0.png}
%\end{minipage}
%\begin{minipage}[b]{0.134\linewidth}
%\includegraphics[width=\linewidth]{1002212612_beta_20.0.png}
%\end{minipage}
%
%\begin{minipage}[b]{0.134\linewidth}
%\includegraphics[width=\linewidth]{1015338126_beta_-14.5.png}
%\end{minipage}
%\begin{minipage}[b]{0.134\linewidth}
%\includegraphics[width=\linewidth]{1015338126_beta_-10.4.png}
%\end{minipage}
%\begin{minipage}[b]{0.134\linewidth}
%\includegraphics[width=\linewidth]{1015338126_beta_-5.0.png}
%\end{minipage}
%\begin{minipage}[b]{0.134\linewidth}
%\includegraphics[width=\linewidth]{1015338126_beta_0.0.png}
%\end{minipage}
%\begin{minipage}[b]{0.134\linewidth}
%\includegraphics[width=\linewidth]{1015338126_beta_5.0.png}
%\end{minipage}
%\begin{minipage}[b]{0.134\linewidth}
%\includegraphics[width=\linewidth]{1015338126_beta_10.0.png}
%\end{minipage}
%\begin{minipage}[b]{0.134\linewidth}
%\includegraphics[width=\linewidth]{1015338126_beta_20.0.png}
%\end{minipage}
%
%\begin{minipage}[b]{0.134\linewidth}
%\includegraphics[width=\linewidth]{1015466717_beta_-20.0.png}
%\end{minipage}
%\begin{minipage}[b]{0.134\linewidth}
%\includegraphics[width=\linewidth]{1015466717_beta_-10.0.png}
%\end{minipage}
%\begin{minipage}[b]{0.134\linewidth}
%\includegraphics[width=\linewidth]{1015466717_beta_-5.0.png}
%\end{minipage}
%\begin{minipage}[b]{0.134\linewidth}
%\includegraphics[width=\linewidth]{1015466717_beta_0.0.png}
%\end{minipage}
%\begin{minipage}[b]{0.134\linewidth}
%\includegraphics[width=\linewidth]{1015466717_beta_5.0.png}
%\end{minipage}
%\begin{minipage}[b]{0.134\linewidth}
%\includegraphics[width=\linewidth]{1015466717_beta_10.0.png}
%\end{minipage}
%\begin{minipage}[b]{0.134\linewidth}
%\includegraphics[width=\linewidth]{1015466717_beta_20.0.png}
%\end{minipage}
%
%\begin{minipage}[b]{0.134\linewidth}
%\includegraphics[width=\linewidth]{1015511213_beta_-20.0.png}
%\end{minipage}
%\begin{minipage}[b]{0.134\linewidth}
%\includegraphics[width=\linewidth]{1015511213_beta_-10.0.png}
%\end{minipage}
%\begin{minipage}[b]{0.134\linewidth}
%\includegraphics[width=\linewidth]{1015511213_beta_-5.0.png}
%\end{minipage}
%\begin{minipage}[b]{0.134\linewidth}
%\includegraphics[width=\linewidth]{1015511213_beta_0.0.png}
%\end{minipage}
%\begin{minipage}[b]{0.134\linewidth}
%\includegraphics[width=\linewidth]{1015511213_beta_5.0.png}
%\end{minipage}
%\begin{minipage}[b]{0.134\linewidth}
%\includegraphics[width=\linewidth]{1015511213_beta_10.0.png}
%\end{minipage}
%\begin{minipage}[b]{0.134\linewidth}
%\includegraphics[width=\linewidth]{1015511213_beta_20.0.png}
%\end{minipage}
%
%\begin{minipage}[b]{0.134\linewidth}
%\includegraphics[width=\linewidth]{1015955655_beta_-20.png}
%\end{minipage}
%\begin{minipage}[b]{0.134\linewidth}
%\includegraphics[width=\linewidth]{1015955655_beta_-10.0.png}
%\end{minipage}
%\begin{minipage}[b]{0.134\linewidth}
%\includegraphics[width=\linewidth]{1015955655_beta_-5.0.png}
%\end{minipage}
%\begin{minipage}[b]{0.134\linewidth}
%\includegraphics[width=\linewidth]{1015955655_beta_0.0.png}
%\end{minipage}
%\begin{minipage}[b]{0.134\linewidth}
%\includegraphics[width=\linewidth]{1015955655_beta_5.0.png}
%\end{minipage}
%\begin{minipage}[b]{0.134\linewidth}
%\includegraphics[width=\linewidth]{1015955655_beta_10.0.png}
%\end{minipage}
%\begin{minipage}[b]{0.134\linewidth}
%\includegraphics[width=\linewidth]{1015955655_beta_20.0.png}
%\end{minipage}
%
%\begin{minipage}[b]{0.134\linewidth}
%\includegraphics[width=\linewidth]{1016429852_beta_-20.0.png}
%\end{minipage}
%\begin{minipage}[b]{0.134\linewidth}
%\includegraphics[width=\linewidth]{1016429852_beta_-10.0.png}
%\end{minipage}
%\begin{minipage}[b]{0.134\linewidth}
%\includegraphics[width=\linewidth]{1016429852_beta_-5.0.png}
%\end{minipage}
%\begin{minipage}[b]{0.134\linewidth}
%\includegraphics[width=\linewidth]{1016429852_beta_0.0.png}
%\end{minipage}
%\begin{minipage}[b]{0.134\linewidth}
%\includegraphics[width=\linewidth]{1016429852_beta_5.0.png}
%\end{minipage}
%\begin{minipage}[b]{0.134\linewidth}
%\includegraphics[width=\linewidth]{1016429852_beta_10.0.png}
%\end{minipage}
%\begin{minipage}[b]{0.134\linewidth}
%\includegraphics[width=\linewidth]{1016429852_beta_20.0.png}
%\end{minipage}
%
%
%% \caption{Manipulating the health attribute with different $\beta$ values for regenerated patients' images. The synthesized image without manipulation is at $\beta=0$ (middle column). These are projections that do not allow identification of real persons}
%
%%\label{fig:manipulations_1}
%
%\end{figure}
%
%% new figure ##############################################################
%
%
%\begin{figure}[h]
%\centering
%\begin{minipage}[b]{0.134\linewidth}
%\centering
%$\beta = -20$
%\end{minipage}
%\begin{minipage}[b]{0.134\linewidth}
%\centering
%$\beta = -10$
%\end{minipage}
%\begin{minipage}[b]{0.134\linewidth}
%\centering
%$\beta = -5$
%\end{minipage}
%\begin{minipage}[b]{0.134\linewidth}
%\centering
%$\xfake, \beta = 0$
%\end{minipage}
%\begin{minipage}[b]{0.134\linewidth}
%\centering
%$\beta = 5$
%\end{minipage}
%\begin{minipage}[b]{0.134\linewidth}
%\centering
%$\beta = 10$ 
%\end{minipage}
%\begin{minipage}[b]{0.134\linewidth}
%\centering
%$\beta = 20$ 
%\end{minipage}
%
%\begin{minipage}[b]{0.134\linewidth}
%\includegraphics[width=\linewidth]{8000282700_beta_-20.0.png}
%\end{minipage}
%\begin{minipage}[b]{0.134\linewidth}
%\includegraphics[width=\linewidth]{8000282700_beta_-10.0.png}
%\end{minipage}
%\begin{minipage}[b]{0.134\linewidth}
%\includegraphics[width=\linewidth]{8000282700_beta_-5.0.png}
%\end{minipage}
%\begin{minipage}[b]{0.134\linewidth}
%\includegraphics[width=\linewidth]{8000282700_beta_0.0.png}
%\end{minipage}
%\begin{minipage}[b]{0.134\linewidth}
%\includegraphics[width=\linewidth]{8000282700_beta_5.0.png}
%\end{minipage}
%\begin{minipage}[b]{0.134\linewidth}
%\includegraphics[width=\linewidth]{8000282700_beta_10.0.png}
%\end{minipage}
%\begin{minipage}[b]{0.134\linewidth}
%\includegraphics[width=\linewidth]{8000282700_beta_20.0.png}
%\end{minipage}
%
%\begin{minipage}[b]{0.134\linewidth}
%\includegraphics[width=\linewidth]{8000075673_beta_-20.0.png}
%\end{minipage}
%\begin{minipage}[b]{0.134\linewidth}
%\includegraphics[width=\linewidth]{8000075673_beta_-10.0.png}
%\end{minipage}
%\begin{minipage}[b]{0.134\linewidth}
%\includegraphics[width=\linewidth]{8000075673_beta_-5.0.png}
%\end{minipage}
%\begin{minipage}[b]{0.134\linewidth}
%\includegraphics[width=\linewidth]{8000075673_beta_0.0.png}
%\end{minipage}
%\begin{minipage}[b]{0.134\linewidth}
%\includegraphics[width=\linewidth]{8000075673_beta_5.0.png}
%\end{minipage}
%\begin{minipage}[b]{0.134\linewidth}
%\includegraphics[width=\linewidth]{8000075673_beta_10.0.png}
%\end{minipage}
%\begin{minipage}[b]{0.134\linewidth}
%\includegraphics[width=\linewidth]{8000075673_beta_20.0.png}
%\end{minipage}
%
%\begin{minipage}[b]{0.134\linewidth}
%\includegraphics[width=\linewidth]{8000316720_beta_-20.0.png}
%\end{minipage}
%\begin{minipage}[b]{0.134\linewidth}
%\includegraphics[width=\linewidth]{8000316720_beta_-10.0.png}
%\end{minipage}
%\begin{minipage}[b]{0.134\linewidth}
%\includegraphics[width=\linewidth]{8000316720_beta_-5.0.png}
%\end{minipage}
%\begin{minipage}[b]{0.134\linewidth}
%\includegraphics[width=\linewidth]{8000316720_beta_0.0.png}
%\end{minipage}
%\begin{minipage}[b]{0.134\linewidth}
%\includegraphics[width=\linewidth]{8000316720_beta_5.0.png}
%\end{minipage}
%\begin{minipage}[b]{0.134\linewidth}
%\includegraphics[width=\linewidth]{8000316720_beta_10.0.png}
%\end{minipage}
%\begin{minipage}[b]{0.134\linewidth}
%\includegraphics[width=\linewidth]{8000316720_beta_20.0.png}
%\end{minipage}
%
%\begin{minipage}[b]{0.134\linewidth}
%\includegraphics[width=\linewidth]{8000341177_beta_-20.0.png}
%\end{minipage}
%\begin{minipage}[b]{0.134\linewidth}
%\includegraphics[width=\linewidth]{8000341177_beta_-10.0.png}
%\end{minipage}
%\begin{minipage}[b]{0.134\linewidth}
%\includegraphics[width=\linewidth]{8000341177_beta_-5.0.png}
%\end{minipage}
%\begin{minipage}[b]{0.134\linewidth}
%\includegraphics[width=\linewidth]{8000341177_beta_0.0.png}
%\end{minipage}
%\begin{minipage}[b]{0.134\linewidth}
%\includegraphics[width=\linewidth]{8000341177_beta_5.0.png}
%\end{minipage}
%\begin{minipage}[b]{0.134\linewidth}
%\includegraphics[width=\linewidth]{8000341177_beta_10.0.png}
%\end{minipage}
%\begin{minipage}[b]{0.134\linewidth}
%\includegraphics[width=\linewidth]{8000341177_beta_20.0.png}
%\end{minipage}
%
%\begin{minipage}[b]{0.134\linewidth}
%\includegraphics[width=\linewidth]{8000785684_beta_-20.0.png}
%\end{minipage}
%\begin{minipage}[b]{0.134\linewidth}
%\includegraphics[width=\linewidth]{8000785684_beta_-10.0.png}
%\end{minipage}
%\begin{minipage}[b]{0.134\linewidth}
%\includegraphics[width=\linewidth]{8000785684_beta_-5.0.png}
%\end{minipage}
%\begin{minipage}[b]{0.134\linewidth}
%\includegraphics[width=\linewidth]{8000785684_beta_0.0.png}
%\end{minipage}
%\begin{minipage}[b]{0.134\linewidth}
%\includegraphics[width=\linewidth]{8000785684_beta_5.0.png}
%\end{minipage}
%\begin{minipage}[b]{0.134\linewidth}
%\includegraphics[width=\linewidth]{8000785684_beta_10.0.png}
%\end{minipage}
%\begin{minipage}[b]{0.134\linewidth}
%\includegraphics[width=\linewidth]{8000785684_beta_20.0.png}
%\end{minipage}
%
%\begin{minipage}[b]{0.134\linewidth}
%\includegraphics[width=\linewidth]{1016059162_beta_-20.0.png}
%\end{minipage}
%\begin{minipage}[b]{0.134\linewidth}
%\includegraphics[width=\linewidth]{1016059162_beta_-10.0.png}
%\end{minipage}
%\begin{minipage}[b]{0.134\linewidth}
%\includegraphics[width=\linewidth]{1016059162_beta_-5.0.png}
%\end{minipage}
%\begin{minipage}[b]{0.134\linewidth}
%\includegraphics[width=\linewidth]{1016059162_beta_0.0.png}
%\end{minipage}
%\begin{minipage}[b]{0.134\linewidth}
%\includegraphics[width=\linewidth]{1016059162_beta_5.0.png}
%\end{minipage}
%\begin{minipage}[b]{0.134\linewidth}
%\includegraphics[width=\linewidth]{1016059162_beta_10.0.png}
%\end{minipage}
%\begin{minipage}[b]{0.134\linewidth}
%\includegraphics[width=\linewidth]{1016059162_beta_20.0.png}
%\end{minipage}
%
%\begin{minipage}[b]{0.134\linewidth}
%\includegraphics[width=\linewidth]{1016080358_beta_-20.0.png}
%\end{minipage}
%\begin{minipage}[b]{0.134\linewidth}
%\includegraphics[width=\linewidth]{1016080358_beta_-10.0.png}
%\end{minipage}
%\begin{minipage}[b]{0.134\linewidth}
%\includegraphics[width=\linewidth]{1016080358_beta_-5.0.png}
%\end{minipage}
%\begin{minipage}[b]{0.134\linewidth}
%\includegraphics[width=\linewidth]{1016080358_beta_0.0.png}
%\end{minipage}
%\begin{minipage}[b]{0.134\linewidth}
%\includegraphics[width=\linewidth]{1016080358_beta_5.0.png}
%\end{minipage}
%\begin{minipage}[b]{0.134\linewidth}
%\includegraphics[width=\linewidth]{1016080358_beta_10.0.png}
%\end{minipage}
%\begin{minipage}[b]{0.134\linewidth}
%\includegraphics[width=\linewidth]{1016080358_beta_20.0.png}
%\end{minipage}
%
%\begin{minipage}[b]{0.134\linewidth}
%\includegraphics[width=\linewidth]{8000955465_beta_-20.0.png}
%\end{minipage}
%\begin{minipage}[b]{0.134\linewidth}
%\includegraphics[width=\linewidth]{8000955465_beta_-10.0.png}
%\end{minipage}
%\begin{minipage}[b]{0.134\linewidth}
%\includegraphics[width=\linewidth]{8000955465_beta_-5.0.png}
%\end{minipage}
%\begin{minipage}[b]{0.134\linewidth}
%\includegraphics[width=\linewidth]{8000955465_beta_0.0.png}
%\end{minipage}
%\begin{minipage}[b]{0.134\linewidth}
%\includegraphics[width=\linewidth]{8000955465_beta_5.0.png}
%\end{minipage}
%\begin{minipage}[b]{0.134\linewidth}
%\includegraphics[width=\linewidth]{8000955465_beta_10.0.png}
%\end{minipage}
%\begin{minipage}[b]{0.134\linewidth}
%\includegraphics[width=\linewidth]{8000955465_beta_20.0.png}
%\end{minipage}
%
%\begin{minipage}[b]{0.134\linewidth}
%\includegraphics[width=\linewidth]{8001006450beta_-20.0.png}
%\end{minipage}
%\begin{minipage}[b]{0.134\linewidth}
%\includegraphics[width=\linewidth]{8001006450beta_-10.0.png}
%\end{minipage}
%\begin{minipage}[b]{0.134\linewidth}
%\includegraphics[width=\linewidth]{8001006450beta_-5.0.png}
%\end{minipage}
%\begin{minipage}[b]{0.134\linewidth}
%\includegraphics[width=\linewidth]{8001006450beta_0.0.png}
%\end{minipage}
%\begin{minipage}[b]{0.134\linewidth}
%\includegraphics[width=\linewidth]{8001006450beta_5.0.png}
%\end{minipage}
%\begin{minipage}[b]{0.134\linewidth}
%\includegraphics[width=\linewidth]{8001006450beta_10.0.png}
%\end{minipage}
%\begin{minipage}[b]{0.134\linewidth}
%\includegraphics[width=\linewidth]{8001006450beta_20.0.png}
%\end{minipage}
%
%\begin{minipage}[b]{0.134\linewidth}
%\includegraphics[width=\linewidth]{8000172176_beta_-20.0.png}
%\end{minipage}
%\begin{minipage}[b]{0.134\linewidth}
%\includegraphics[width=\linewidth]{8000172176_beta_-10.0.png}
%\end{minipage}
%\begin{minipage}[b]{0.134\linewidth}
%\includegraphics[width=\linewidth]{8000172176_beta_-5.0.png}
%\end{minipage}
%\begin{minipage}[b]{0.134\linewidth}
%\includegraphics[width=\linewidth]{8000172176_beta_0.0.png}
%\end{minipage}
%\begin{minipage}[b]{0.134\linewidth}
%\includegraphics[width=\linewidth]{8000172176_beta_5.0.png}
%\end{minipage}
%\begin{minipage}[b]{0.134\linewidth}
%\includegraphics[width=\linewidth]{8000172176_beta_10.0.png}
%\end{minipage}
%\begin{minipage}[b]{0.134\linewidth}
%\includegraphics[width=\linewidth]{8000172176_beta_20.0.png}
%\end{minipage}
%
%\begin{minipage}[b]{0.134\linewidth}
%\includegraphics[width=\linewidth]{1002154510_beta_-20.0.png}
%\end{minipage}
%\begin{minipage}[b]{0.134\linewidth}
%\includegraphics[width=\linewidth]{1002154510_beta_-10.0.png}
%\end{minipage}
%\begin{minipage}[b]{0.134\linewidth}
%\includegraphics[width=\linewidth]{1002154510_beta_-5.0.png}
%\end{minipage}
%\begin{minipage}[b]{0.134\linewidth}
%\includegraphics[width=\linewidth]{1002154510_beta_0.0.png}
%\end{minipage}
%\begin{minipage}[b]{0.134\linewidth}
%\includegraphics[width=\linewidth]{1002154510_beta_5.0.png}
%\end{minipage}
%\begin{minipage}[b]{0.134\linewidth}
%\includegraphics[width=\linewidth]{1002154510_beta_10.0.png}
%\end{minipage}
%\begin{minipage}[b]{0.134\linewidth}
%\includegraphics[width=\linewidth]{1002154510_beta_20.0.png}
%\end{minipage}
%
%\caption{Manipulating the health attribute with different $\beta$ values for regenerated patients' images. The synthesized image without manipulation is at $\beta=0$ (middle column). These are projections that do not allow identification of real persons}
%
%\label{fig:manipulations_2}
%\end{figure}

% \TODO{*New in 2023* Data use declaration and acknowledgment: Authors must declare the data origin, data license, and (when appropriate) ethics application number}
%\bibliographystyle{splncs04}
%\bibliography{ref}
\begin{thebibliography}{10}
\providecommand{\url}[1]{\texttt{#1}}
\providecommand{\urlprefix}{URL }
\providecommand{\doi}[1]{https://doi.org/#1}

\bibitem{amin2017eighth}
Amin, M.B., Greene, F.L., Edge, S.B., Compton, C.C., Gershenwald, J.E.,
  Brookland, R.K., Meyer, L., Gress, D.M., Byrd, D.R., Winchester, D.P.: The
  eighth edition ajcc cancer staging manual: continuing to build a bridge from
  a population-based to a more “personalized” approach to cancer staging.
  CA: a cancer journal for clinicians  \textbf{67}(2),  93--99 (2017)

\bibitem{bagnis2020judging}
Bagnis, A., Caffo, E., Cipolli, C., De~Palma, A., Farina, G., Mattarozzi, K.:
  Judging health care priority in emergency situations: patient facial
  appearance matters. Social Science \& Medicine  \textbf{260},  113180 (2020)

\bibitem{chen2022mri}
Chen, J., Lu, S., Mao, Y., Tan, L., Li, G., Gao, Y., Tan, P., Huang, D., Zhang,
  X., Qiu, Y., et~al.: An mri-based radiomics-clinical nomogram for the overall
  survival prediction in patients with hypopharyngeal squamous cell carcinoma:
  a multi-cohort study. European Radiology pp. 1--10 (2022)

\bibitem{emura2012survival}
Emura, T., Chen, Y.H., Chen, H.Y.: Survival prediction based on compound
  covariate under cox proportional hazard models. PLoS One  \textbf{7}(10),
  e47627 (2012)

\bibitem{gurovich2019identifying}
Gurovich, Y., Hanani, Y., Bar, O., Nadav, G., Fleischer, N., Gelbman, D.,
  Basel-Salmon, L., Krawitz, P.M., Kamphausen, S.B., Zenker, M., et~al.:
  Identifying facial phenotypes of genetic disorders using deep learning.
  Nature medicine  \textbf{25}(1),  60--64 (2019)

\bibitem{huang2022deep}
Huang, Y., Bert, C., Sommer, P., Frey, B., Gaipl, U., Distel, L.V., Weissmann,
  T., Uder, M., Schmidt, M.A., D{\"o}rfler, A., et~al.: Deep learning for brain
  metastasis detection and segmentation in longitudinal mri data. Medical
  Physics  \textbf{49}(9),  5773--5786 (2022)

\bibitem{huang2023benchmarking}
Huang, Y., Gomaa, A., Weissmann, T., Grigo, J., Tkhayat, H.B., Frey, B., Gaipl,
  U.S., Distel, L.V., Maier, A., Fietkau, R., et~al.: Benchmarking chatgpt-4 on
  acr radiation oncology in-training ({TXIT}) exam and red journal gray zone
  cases: Potentials and challenges for ai-assisted medical education and
  decision making in radiation oncology. arXiv preprint  (2023)

\bibitem{hui2015diagnostic}
Hui, D., Hess, K., Santos, R.d., Chisholm, G., Bruera, E.: A diagnostic model
  for impending death in cancer patients: preliminary report. Cancer
  \textbf{121}(21),  3914--3921 (2015)

\bibitem{karras2020training}
Karras, T., Aittala, M., Hellsten, J., Laine, S., Lehtinen, J., Aila, T.:
  Training generative adversarial networks with limited data. Advances in
  neural information processing systems  \textbf{33},  12104--12114 (2020)

\bibitem{karras2019style}
Karras, T., Laine, S., Aila, T.: A style-based generator architecture for
  generative adversarial networks. In: Proc. CVPR. pp. 4401--4410 (2019)

\bibitem{karras2020analyzing}
Karras, T., Laine, S., Aittala, M., Hellsten, J., Lehtinen, J., Aila, T.:
  Analyzing and improving the image quality of stylegan. In: Proc. CVPR. pp.
  8110--8119 (2020)

\bibitem{katzman2018deepsurv}
Katzman, J.L., Shaham, U., Cloninger, A., Bates, J., Jiang, T., Kluger, Y.:
  Deepsurv: personalized treatment recommender system using a cox proportional
  hazards deep neural network. BMC medical research methodology
  \textbf{18}(1),  1--12 (2018)

\bibitem{kidd2018survival}
Kidd, A.C., McGettrick, M., Tsim, S., Halligan, D.L., Bylesjo, M., Blyth, K.G.:
  Survival prediction in mesothelioma using a scalable lasso regression model:
  instructions for use and initial performance using clinical predictors. BMJ
  open respiratory research  \textbf{5}(1),  e000240 (2018)

\bibitem{kim2019deep}
Kim, D.W., Lee, S., Kwon, S., Nam, W., Cha, I.H., Kim, H.J.: Deep
  learning-based survival prediction of oral cancer patients. Scientific
  reports  \textbf{9}(1),  1--10 (2019)

\bibitem{dlib09}
King, D.E.: Dlib-ml: A machine learning toolkit. Journal of Machine Learning
  Research  \textbf{10},  1755--1758 (2009)

\bibitem{kong2018automatic}
Kong, X., Gong, S., Su, L., Howard, N., Kong, Y.: Automatic detection of
  acromegaly from facial photographs using machine learning methods.
  EBioMedicine  \textbf{27},  94--102 (2018)

\bibitem{liang2020identification}
Liang, B., Yang, N., He, G., Huang, P., Yang, Y.: Identification of the facial
  features of patients with cancer: a deep learning--based pilot study. Journal
  of Medical Internet Research  \textbf{22}(4),  e17234 (2020)

\bibitem{lin2002modeling}
Lin, H., Zelterman, D.: Modeling survival data: extending the cox model (2002)

\bibitem{oken1982toxicity}
Oken, M.M., Creech, R.H., Tormey, D.C., Horton, J., Davis, T.E., McFadden,
  E.T., Carbone, P.P.: Toxicity and response criteria of the eastern
  cooperative oncology group. Am. J. Clin. Oncol.  \textbf{5}(6),  649--656
  (1982)

\bibitem{de2017quality}
de~Oliveira, W.A.: Quality of life, facial appearance and self-esteem in
  patients with orthodontic treatment. Revista Mexicana de Ortodoncia
  \textbf{5}(3),  138--139 (2017)

\bibitem{putz2023segment}
Putz, F., Grigo, J., Weissmann, T., Schubert, P., Hoefler, D., Gomaa, A.,
  Tkhayat, H.B., Hagag, A., Lettmaier, S., Frey, B., et~al.: The segment
  anything foundation model achieves favorable brain tumor autosegmentation
  accuracy on mri to support radiotherapy treatment planning. arXiv preprint
  (2023)

\bibitem{rankin2003perceived}
Rankin, M., Borah, G.L.: Perceived functional impact of abnormal facial
  appearance. Plastic and reconstructive surgery  \textbf{111}(7),  2140--2146
  (2003)

\bibitem{su2021deep}
Su, Z., Liang, B., Shi, F., Gelfond, J., {\v{S}}egalo, S., Wang, J., Jia, P.,
  Hao, X.: Deep learning-based facial image analysis in medical research: a
  systematic review protocol. BMJ open  \textbf{11}(11),  e047549 (2021)

\bibitem{van2017survival}
van Timmeren, J.E., Leijenaar, R.T., van Elmpt, W., Reymen, B., Oberije, C.,
  Monshouwer, R., Bussink, J., Brink, C., Hansen, O., Lambin, P.: Survival
  prediction of non-small cell lung cancer patients using radiomics analyses of
  cone-beam ct images. Radiotherapy and Oncology  \textbf{123}(3),  363--369
  (2017)

\bibitem{vale2021long}
Vale-Silva, L.A., Rohr, K.: Long-term cancer survival prediction using
  multimodal deep learning. Scientific Reports  \textbf{11}(1),  1--12 (2021)

\bibitem{wallis2022clever}
Wallis, D., Buvat, I.: Clever hans effect found in a widely used brain tumour
  mri dataset. Medical Image Analysis  \textbf{77},  102368 (2022)

\bibitem{wang2021cox}
Wang, D., Jing, Z., He, K., Garmire, L.X.: Cox-nnet v2. 0: improved
  neural-network-based survival prediction extended to large-scale emr data.
  Bioinformatics  \textbf{37}(17),  2772--2774 (2021)

\bibitem{wankhede2022dynamic}
Wankhede, D.S., Rangasamy, S.: Dynamic architecture based deep learning
  approach for glioblastoma brain tumor survival prediction. Neuroscience
  Informatics p. 100062 (2022)

\bibitem{weissmann2023deep}
Weissmann, T., Huang, Y., Fischer, S., Roesch, J., Mansoorian, S., Ayala~Gaona,
  H., Gostian, A.O., Hecht, M., Lettmaier, S., Deloch, L., et~al.: Deep
  learning for automatic head and neck lymph node level delineation provides
  expert-level accuracy. Frontiers in Oncology  \textbf{13},  1115258 (2023)

\bibitem{withington2005hippocratic}
Withington, E., Lonie, I., Chadwick, J., Mann, W.N., Lloyd, G., et~al.:
  Hippocratic writings. Penguin UK (2005)

\bibitem{yolcu2017deep}
Yolcu, G., Oztel, I., Kazan, S., Oz, C., Palaniappan, K., Lever, T.E., Bunyak,
  F.: Deep learning-based facial expression recognition for monitoring
  neurological disorders. In: 2017 IEEE International Conference on
  Bioinformatics and Biomedicine (BIBM). pp. 1652--1657. IEEE (2017)

\bibitem{zhu2016deep}
Zhu, X., Yao, J., Huang, J.: Deep convolutional neural network for survival
  analysis with pathological images. In: 2016 IEEE International Conference on
  Bioinformatics and Biomedicine (BIBM). pp. 544--547. IEEE (2016)

\end{thebibliography}


\end{document}
